\documentclass[10pt,a4paper]{amsart}
\usepackage[margin=19mm]{geometry}
\usepackage{amsmath,amssymb,mathtools}
\usepackage[scr=rsfs,cal=euler]{mathalfa}
\usepackage{xfrac}
\usepackage[unicode,hidelinks,bookmarks=false]{hyperref}
\newcommand{\ie}{i.e.\ }
\newcommand{\cf}{cf.\ }
\newcommand{\resp}{respectively\ }
\newcommand{\Subsection}{\subsection}
\providecommand{\nobreakdash}{\nobreak\hbox{-}}
\setlength{\emergencystretch}{2em}
\multlinegap0pt
\usepackage{amssymb}
\usepackage{multirow}
\usepackage{stackengine}
\usepackage{tcolorbox}
\newtheorem{proposition}{Proposition}
\newcommand{\vol}{\mathop{\rm vol}\nolimits}
\title{Variational formulations of transport phenomena on combinatorial meshes}
\author{Kiprian Berbatov, Andrey P. Jivkov}
\date{Source: arXiv:2505.09443v8 (CC BY 4.0)}
\begin{document}
\maketitle

\noindent\emph{Source and licence.} Variational formulations of transport phenomena on combinatorial meshes. Version arXiv:2505.09443v8 (CC BY 4.0), distributed by arXiv under the Creative Commons Attribution 4.0 International licence, http://creativecommons.org/licenses/by/4.0/, as recorded in the arXiv licence metadata for this exact version and shown on the abstract page https://arxiv.org/abs/2505.09443v8. Full licence notice: \texttt{licenses/arxiv-2505.09443-CC-BY-4.0.txt}.\par\smallskip

\noindent\textit{Editorial scope.} Counted unit: the article's proposition proposition (source0.tex lines 833--839). The article's complete original proof of it is delivered with it (source0.tex lines 840--861, one \texttt{proof} environment of 22 lines). No mathematics is written, repaired or re-derived: the statement and the proof are the article's own lines, in the article's own notation, and no retained line is substituted or re-flowed.  No further source-local context is needed: every cross-reference of every retained line resolves inside the two delivered spans.  Nothing was omitted from any retained span, so the package carries no omission record.  No retained line contains a citation, so the wrapper carries no bibliography and no bibliography entry.  No retained line contains a figure environment, an image-inclusion command or drawing code, so no image is delivered and the package declares no asset dependency.  Editorial formatting only, in addition: an AMS \texttt{amsart} preamble that restates the article's own declarations and macros used by the retained lines, namely the environments \texttt{proposition} on separate counters, together with the standard packages those lines need.  The retained environments print in this document as Proposition 1 in order of appearance, whereas the article prints the same items with the numbering of its own full document.  The complete native source of the article as distributed in the arXiv e-print for this exact version is bundled unchanged as \texttt{sources/178116/source0.tex}.
\begin{proposition}
  Let $(M, g)$ be an oriented Riemannian manifold of dimension $D$ with a volume form $\vol$, $0 \leq p \leq D$.
  Then
  \begin{equation}
    \star_{D - p} \circ \star_p = (-1)^{p (D - p)} {\rm id}_{\Omega^p M}.
  \end{equation}
\end{proposition}

\begin{proof}
  Consider an oriented orthonormal basis $(e^1, ..., e^D)$ of $\Omega^1 M$, $\vol := e^1 \wedge ... \wedge e^D$ be the volume form.
  Let $I = (I_1, ..., I_p)$, $1 \leq I_1 < ... < I_p \leq D$ and $e_I := e^{I_1} \wedge ... \wedge e^{I_p}$.
  Let $J$ be the ordered complement of $I$ (with respect to $(1, ..., D)$),
  and $e^J := e^{J_1} \wedge ... \wedge e^{J_{D - p}}$.
  Then $\star_p e^I = s e^J$ for some $s \in \{-1, 1\}$ (since taking the wedge product gives $e^1 \wedge ... \wedge e^D = \vol$ after reordering, taken into account by $s$), and similarly $\star_{D - p} e^J = t e^I$ for some $t \in \{-1, 1\}$.
  Then
  \begin{equation}
    s\, \vol
    = g^*_{D - p}(\star_p e^I, e^J)\, \vol = e^I \wedge e^J
    = (-1)^{p (D - p)} e^J \wedge e^I
    = (-1)^{p (D - p)} g^*_p(\star_{D - p} e^J, e^I)\, \vol
    = (-1)^{p (D - p)} t\, \vol,
  \end{equation}
  from which it follows that $s = (-1)^{p (D - p)} t$ or $s t = (-1)^{p (D - p)}$.
  Hence,
  \begin{equation}
    \star_{D - p} \star_p e^I = s t e^I = (-1)^{p (D - p)} e^I.
  \end{equation}
  Since $\star$ is $(\mathcal{F} M)$-linear, it follows that for any $\omega \in \Omega^p M$,
  $\star_{D - p} \star_p \omega = (-1)^{p (D - p)} \omega$.
\end{proof}

\end{document}
